% ************************** Thesis Abstract *****************************
% Use `abstract' as an option in the document class to print only the titlepage and the abstract.
\begin{abstract}
Background: 
The prevalence of stress and burnout in healthcare professionals due to excessive administrative work is a well-known problem since the advent of electronic health records (EHR) in 2016. Especially nursing personnel has been found to spend up to 60\% of their time for handling documentation, e.g. for compiling progress notes for their patients (Khamisa et al, 2013). In Greece, the problem is even worse, as there still exist handwritten patient records that take up considerable time and energy from nurses to write, read, interpret, and most importantly retrieve on demand. The goal of this project is to ``liberate`` nurse staff from time-consuming logistics problems, and offer them the opportunity to focus on patient treatment and care instead. The way to achieve this is by designing a voice chatbot that will automate and expedite the information retrieval process from patient data files. 

Objective:
This study aims to design, develop, and test an artificial intelligence–powered voice assistant to help psychiatric nurses manage the daily interactions with their patients.

Methods:
The voice assistant is being implemented in several steps. First, a comprehensive mock knowledge base is created, which defines real-world user profiles. Second, conversation behavior between users and the voice assistant is defined via a well-designed LLM prompt. Third, advanced features like memory, tools, and an orchestrating ``agent`` are implemented to provide continuity, reasoning, and autonomy. Fourth, two separate system evaluations are conducted via use cases (scenarios) using RAGAs framework as well as additional qualitative evaluations. 

Results:
A prototype voice assistant was evaluated in using various use cases. Preliminary qualitative test results demonstrate reasonable rates of dialogue success and answer correctness.

Conclusions:
This study suggests the feasibility of using an intelligent voice assistant to help nurses in managing patients in psychiatric clinics.
\end{abstract}
